\chapter{Overfitting e regolarizzazione}
\label{cap:regolarizzazione}

\begin{inbreve}
Un modello che si adatta perfettamente ai dati di training può predire male i dati nuovi. Il capitolo definisce
l'overfitting e lo osserva in un caso di studio: il fitting di polinomi, che è una espansione in basi lineare.
Introduce poi lo strumento principale per controllare la complessità di un modello, la regolarizzazione di
Tikhonov (\emph{ridge regression}, \emph{weight decay}), con le varianti $L^1$ ed \emph{elastic net}.
\end{inbreve}

\section{Generalizzazione e overfitting}
\label{sec:overfitting}

\begin{definizione}[Ipotesi dell'apprendimento induttivo]
\emph{Qualsiasi ipotesi $h$ che approssima bene $f$ sugli esempi di training approssimerà bene $f$ anche sulle
istanze nuove, mai viste.}
\end{definizione}

\uline{Questa ipotesi non è sempre vera}, ed è proprio qui che nasce l'overfitting.

\begin{definizione}[Overfitting]
Un learner fa \termine{overfitting} se restituisce un'ipotesi $h\in H$ con errore vero (di generalizzazione,
\emph{rischio}) $R$ ed errore empirico (di training) $E$, quando esiste un'altra $h'\in H$ con
\[
E' > E \qquad\text{ma}\qquad R' < R :
\]
$h'$ si adatta \emph{peggio} ai dati di training ma è \emph{migliore} sui dati nuovi. L'\termine{underfitting} è
la situazione opposta: il modello è troppo semplice per cogliere l'andamento della funzione target, e
l'errore è alto già sul training.
\end{definizione}

Il punto critico diventa quindi la \textbf{stima dell'accuratezza}, che si affronta in modo \textbf{teorico}
(Statistical Learning Theory, capitolo~\ref{cap:slt}) o \textbf{empirico} (errori su training e test,
cross-validation, capitolo~\ref{cap:validazione}).

\section{Un caso di studio: il fitting polinomiale}
\label{sec:polinomi}

Si usa un problema di regressione con un modello parametrico: lo spazio delle ipotesi sono i
\textbf{polinomi di grado $M$} (una LBE con $\phi_j(x) = x^j$, \rif{sec:lbe})
\[
h_\w(x) = w_0 + w_1x + w_2x^2 + \dots + w_Mx^M = \sum_{j=0}^{M}w_jx^j ,
\]
la cui complessità cresce con $M$. La funzione target è $\sin(2\pi x)$ e i target dei $l=10$ esempi contengono
\textbf{rumore gaussiano}, per cui i campioni non stanno esattamente sulla curva vera. Il task è artificiale e
irrealistico (una sola variabile, funzione target nota) e serve solo a visualizzare i concetti. I parametri si
trovano minimizzando la somma degli errori quadratici
\[
E(\w) = \sum_{p=1}^{l}\big(d_p - h_\w(x_p)\big)^2 ,
\]
con le equazioni normali del \rif{sec:equazioni-normali} applicate alle feature $(1,x,x^2,\dots,x^M)$.

\subsection{Effetto del grado del polinomio}

\begin{figure}[ht]
\centering
\begin{tikzpicture}
\begin{groupplot}[group style={group size=4 by 1, horizontal sep=0.35cm}, scatola, width=4.35cm, height=3.7cm,
  xmin=-0.05, xmax=1.05, ymin=-1.6, ymax=1.6, xtick={0,1}, ytick={-1,0,1}, yticklabels={,,}, xlabel={$x$}]
\nextgroupplot[title={\small $M=0$}, yticklabels={$-1$,$0$,$1$}]
\addplot[thick, verde] table {dati/poly_vera.dat};
\addplot[thick, rosso] table {dati/poly_M0.dat};
\addplot[only marks, mark=o, mark size=1.6pt, blu, thick] table {dati/poly_tr10.dat};
\nextgroupplot[title={\small $M=1$}]
\addplot[thick, verde] table {dati/poly_vera.dat};
\addplot[thick, rosso] table {dati/poly_M1.dat};
\addplot[only marks, mark=o, mark size=1.6pt, blu, thick] table {dati/poly_tr10.dat};
\nextgroupplot[title={\small $M=3$}]
\addplot[thick, verde] table {dati/poly_vera.dat};
\addplot[thick, rosso] table {dati/poly_M3.dat};
\addplot[only marks, mark=o, mark size=1.6pt, blu, thick] table {dati/poly_tr10.dat};
\nextgroupplot[title={\small $M=9$}]
\addplot[thick, verde] table {dati/poly_vera.dat};
\addplot[thick, rosso] table {dati/poly_M9.dat};
\addplot[only marks, mark=o, mark size=1.6pt, blu, thick] table {dati/poly_tr10.dat};
\end{groupplot}
\end{tikzpicture}
\caption{Fitting con polinomi di grado $M = 0,1,3,9$ (rosso) dei 10 campioni rumorosi (blu) della funzione vera
$\sin(2\pi x)$ (verde).}
\label{fig:polinomi}
\end{figure}

\begin{itemize}
  \item \textbf{$M=0$} (costante) e \textbf{$M=1$} (retta): il modello è troppo semplice e non coglie l'andamento
        della funzione target: \textbf{underfitting}.
  \item \textbf{$M=3$}: la flessibilità in più è utile e la curva approssima bene il seno.
  \item \textbf{$M=9$}: con 10 punti e 10 parametri il polinomio passa \textbf{esattamente} per tutti i dati,
        $E(\w)=0$ sul training. Ma oscilla violentemente tra i punti e rappresenta malissimo la funzione vera: il
        modello è troppo complesso e \textbf{impara anche il rumore}: \textbf{overfitting}.
\end{itemize}

\subsection{Errore di training e di test al variare di $M$}

Per confrontare modelli e dataset si usa il \termine{Root-Mean-Square error}
\[
E_{RMS} = \sqrt{\frac1l\sum_{p=1}^{l}\big(d_p - h_{\w^*}(x_p)\big)^2}
\]
dove $\w^*$ sono i pesi del modello addestrato (equivale a $\sqrt{2E(\w^*)/l}$ se $E$ è definito con un fattore
$\tfrac12$ davanti alla somma). La divisione per $l$ permette di confrontare dataset di dimensione diversa, la
radice riporta l'errore sulla scala del target.

\begin{figure}[ht]
\centering
\begin{tikzpicture}
\begin{axis}[stile, width=8cm, height=5cm, xmin=-0.3, xmax=9.4, ymin=0, ymax=1.25, xlabel={$M$}, ylabel={$E_{RMS}$},
  xtick={0,3,6,9}, ytick={0,0.5,1}, yticklabels={$0$,$0{,}5$,$1$}, legend pos=north east]
\addplot[thick, blu, mark=o, mark size=1.6pt] table[x=M, y=train] {dati/poly_rms_M.dat}; \addlegendentry{training}
\addplot[thick, rosso, mark=o, mark size=1.6pt, restrict y to domain=0:1.25] table[x=M, y=test] {dati/poly_rms_M.dat}; \addlegendentry{test}
\end{axis}
\end{tikzpicture}
\caption{$E_{RMS}$ di training e di test (1000 punti nuovi) al variare del grado $M$: a sinistra underfitting, a
destra overfitting (per $M=9$ l'errore di test esce dal grafico).}
\label{fig:rms-M}
\end{figure}

L'errore di training \textbf{decresce sempre} all'aumentare di $M$, perché un modello più flessibile può
adattarsi meglio ai dati. L'errore di test prima decresce, poi resta stabile e infine \textbf{esplode} per
$M=9$: è la firma dell'overfitting. Anche i coefficienti ottimi $\w^*$ sono rivelatori
(Tabella~\ref{tab:coeff-M}, calcolata sugli stessi dati della \figrif{fig:polinomi}).

\begin{table}[ht]
\centering\small
\caption{Coefficienti ottimi $\w^*$ dei polinomi di grado $M$ della \figrif{fig:polinomi}.}
\label{tab:coeff-M}
% generato da script/genera_dati.py: coefficienti dei polinomi della figura (M = 0, 1, 3, 9)
\begin{tabular}{lrrrr}
\toprule
 & $M=0$ & $M=1$ & $M=3$ & $M=9$ \\
\midrule
$w_0^*$ & $0{,}18$ & $0{,}81$ & $0{,}31$ & $0{,}35$ \\
$w_1^*$ &  & $-1{,}27$ & $7{,}94$ & $234{,}52$ \\
$w_2^*$ &  &  & $-25{,}34$ & $-5374{,}81$ \\
$w_3^*$ &  &  & $17{,}33$ & $49046{,}92$ \\
$w_4^*$ &  &  &  & $-233841{,}49$ \\
$w_5^*$ &  &  &  & $645849{,}46$ \\
$w_6^*$ &  &  &  & $-1070953{,}26$ \\
$w_7^*$ &  &  &  & $1050929{,}32$ \\
$w_8^*$ &  &  &  & $-562020{,}21$ \\
$w_9^*$ &  &  &  & $126129{,}45$ \\
\bottomrule
\end{tabular}

\end{table}

\begin{idea}[Pesi enormi = overfitting]
Nel polinomio di grado 9 i coefficienti diventano \textbf{enormi} e di segno alterno. Si compensano a vicenda
per far passare la curva esattamente per ogni punto, e producono oscillazioni fortissime. Penalizzare i pesi
grandi è quindi un modo per controllare la complessità: è la regolarizzazione (\rif{sec:tikhonov}).
\end{idea}

\subsection{Effetto della quantità di dati}

\begin{figure}[ht]
\centering
\begin{tikzpicture}
\begin{groupplot}[group style={group size=2 by 1, horizontal sep=0.5cm}, scatola, width=6cm, height=4.2cm,
  xmin=-0.05, xmax=1.05, ymin=-1.7, ymax=1.7, xtick={0,1}, ytick={-1,0,1}, xlabel={$x$}]
\nextgroupplot[title={\small $M=9$, $l=15$}]
\addplot[thick, verde] table {dati/poly_vera.dat};
\addplot[thick, rosso] table {dati/poly9_l15.dat};
\addplot[only marks, mark=o, mark size=1.5pt, blu] table {dati/poly9_l15_pts.dat};
\nextgroupplot[title={\small $M=9$, $l=100$}, yticklabels={,,}]
\addplot[thick, verde] table {dati/poly_vera.dat};
\addplot[only marks, mark=o, mark size=1.1pt, blu] table {dati/poly9_l100_pts.dat};
\addplot[thick, rosso] table {dati/poly9_l100.dat};
\end{groupplot}
\end{tikzpicture}
\caption{Polinomio di grado 9 addestrato su $l=15$ e $l=100$ esempi: con più dati l'overfitting si riduce.}
\label{fig:poly-dati}
\end{figure}

A parità di modello ($M=9$), con 15 esempi il polinomio oscilla ancora ma molto meno, e con 100 esempi segue
quasi perfettamente il seno (\figrif{fig:poly-dati}). \textbf{Con più dati si possono usare modelli più
complessi} senza cadere nell'overfitting.

\begin{idea}[Cosa mostra l'esempio]
La capacità di generalizzazione (misurata come rischio o errore di test) va studiata in relazione a
\textbf{errore di training}, \textbf{complessità del modello} e \textbf{numero di dati}. In generale, al crescere
della complessità l'errore di training decresce sempre. Quello di test invece decresce e poi risale
(\figrif{fig:curva-complessita}). La teoria che lega i tre elementi è la Statistical Learning Theory
(capitolo~\ref{cap:slt}).
\end{idea}

\begin{figure}[ht]
\centering
\begin{tikzpicture}
\begin{axis}[stile, width=8cm, height=4.8cm, domain=0.3:9.7, xmin=0, xmax=10, ymin=0, ymax=1.1, xtick=\empty, ytick=\empty,
  xlabel={complessità del modello}, ylabel={errore}, legend pos=north east]
\addplot[thick, blu] {0.05+0.9*exp(-x/2.3)}; \addlegendentry{training}
\addplot[thick, rosso] {0.2+0.75*exp(-x/2.0)+0.006*x^2}; \addlegendentry{test}
\draw[densely dashed, grigio] (axis cs:4.6,0) -- (axis cs:4.6,1.0);
\node[font=\footnotesize, grigio] at (axis cs:2.2,0.98) {underfitting};
\node[font=\footnotesize, grigio] at (axis cs:7.4,0.55) {overfitting};
\end{axis}
\end{tikzpicture}
\caption{Andamento tipico (qualitativo) degli errori di training e di test al crescere della complessità del
modello.}
\label{fig:curva-complessita}
\end{figure}

\begin{esempio}[Rumore o complessità?]
Il rumore è la causa dell'overfitting, o lo è il compromesso di complessità? Si può avere overfitting anche con
dati perfettamente puliti. Un polinomio di grado molto alto che interpola pochi punti di una funzione liscia può
oscillare molto tra un punto e l'altro. Il problema è la \textbf{complessità eccessiva rispetto ai dati
disponibili}. Per lo stesso motivo un errore di training nullo non implica una buona soluzione.
\end{esempio}

\section{Regolarizzazione di Tikhonov}
\label{sec:tikhonov}

Tra i molti modi di controllare la complessità c'è il \ingl{coefficient shrinkage}, il ``restringimento'' dei
coefficienti. La \termine{ridge regression} (\termine{regolarizzazione di Tikhonov}) aggiunge all'errore un
termine che \textbf{penalizza i pesi grandi}:
\begin{equation}\label{eq:ridge}
\boxed{\;\mathrm{Loss}(\w) = \underbrace{\sum_{p=1}^{l}\big(d_p - \x_p\T\w\big)^2}_{\text{termine d'errore }(\Remp)}
+ \underbrace{\lambda\,\norm{\w}^2}_{\text{penalità}},\qquad \norm{\w}^2 = \sum_j w_j^2 .\;}
\end{equation}
$\lambda>0$ è l'\termine{iperparametro di regolarizzazione}, scelto in fase di \textbf{model selection}
(capitolo~\ref{cap:validazione}). Il modello risultante è più ``liscio'' (\ingl{smooth}). Si favoriscono
modelli con pesi piccoli o nulli, che usano di fatto meno termini, cioè meno complessi.

\begin{attenzione}[Loss ed errore non sono più sinonimi]
Da qui in poi \textbf{loss} indica la funzione obiettivo usata per l'addestramento (errore più penalità) ed
\textbf{errore} $E$ il solo termine sui dati, che misura la qualità del modello. Nei risultati si riporta
l'errore, perché è la misura utile per chi usa il modello.
\end{attenzione}

\begin{idea}[Cambia l'obiettivo dell'apprendimento]
Non si cerca più \emph{solo} il modello che si adatta meglio ai dati, ma \textbf{il modello più semplice che si
adatta bene ai dati}. La loss è una bilancia tra due piatti:
\begin{itemize}
  \item \textbf{adattamento ai dati} (errore di training): spinge a seguire i dati e minimizzare i residui;
  \item \textbf{semplicità del modello} (penalità $\norm{\w}^2$): punisce i pesi grandi, rendendo il modello più
        liscio e più capace di generalizzare;
\end{itemize}
e $\lambda$ decide quanto pesa la penalità.
\end{idea}

\subsection{Soluzione}

\begin{proprieta}{Soluzione diretta della ridge regression}
\[
\w = (X\T X + \lambda I)^{-1}X\T\dd ,
\]
e per $\lambda>0$ la matrice $X\T X+\lambda I$ è \textbf{sempre invertibile}.
\end{proprieta}
\begin{proof}
In forma matriciale $\mathrm{Loss}(\w) = \norm{\dd - X\w}^2 + \lambda\,\w\T\w$. Il gradiente del primo termine è
$-2X\T(\dd-X\w)$ (\rif{sec:equazioni-normali}); quello del secondo è $2\lambda\w$, perché
$\partial(\lambda\sum_jw_j^2)/\partial w_j = 2\lambda w_j$. Imponendo il gradiente nullo:
\[
-2X\T\dd + 2X\T X\w + 2\lambda\w = \vv0 \iff (X\T X + \lambda I)\w = X\T\dd .
\]
Per ogni $\vv v\neq\vv0$ vale $\vv v\T(X\T X+\lambda I)\vv v = \norm{X\vv v}^2 + \lambda\norm{\vv v}^2 > 0$. La
matrice è quindi definita positiva, e perciò invertibile. È anche un vantaggio numerico rispetto alle equazioni normali
semplici.
\end{proof}

Con la \textbf{discesa del gradiente} si segue ancora il gradiente negativo della loss. Poiché la derivata della
penalità rispetto a $w_j$ è $2\lambda w_j$, la regola diventa
\begin{equation}\label{eq:weight-decay}
\boxed{\;\w_{\mathit{new}} = \w + \eta\,\Delta\w - 2\lambda\w\;}
\end{equation}
dove $\eta\,\Delta\w$ è il passo sul solo termine d'errore. È la tecnica del \termine{weight decay}: anche con
gradiente d'errore nullo, a ogni passo ogni peso si riduce di una frazione del suo valore.

\begin{attenzione}[Il $\lambda$ della regola non è quello della loss]
Un passo di discesa sulla loss completa darebbe $\w_{\mathit{new}} = \w + \eta\,\Delta\w - 2\eta\lambda\w$. Nella
regola \eqref{eq:weight-decay} il fattore $\eta$ della penalità è stato assorbito in $\lambda$: il $\lambda$ della
regola vale $\eta$ volte quello della loss \eqref{eq:ridge}. I due valori numerici quindi non vanno confrontati
direttamente con quello della soluzione diretta.
\end{attenzione}

\subsection{Il compromesso governato da $\lambda$}

\begin{itemize}
  \item \textbf{$\lambda$ piccolo}: conta soprattutto il termine d'errore. Il modello resta troppo complesso
        (pesi di norma alta), con rischio di \textbf{overfitting}.
  \item \textbf{$\lambda$ grande}: conta soprattutto la penalità. L'errore sui dati può crescere troppo, verso
        l'\textbf{underfitting}.
\end{itemize}
Il vantaggio è una \textbf{realizzazione concreta del controllo della complessità}, facile da implementare e di
applicabilità generale: la stessa penalità si usa nelle reti neurali (\rif{sec:reg-reti}). Spingendo i pesi
verso valori piccoli si ottiene un modello con VC-dimension più bassa (\rif{sec:srm}). Il compromesso è regolato
da \textbf{un solo parametro}.

\subsection{Effetto sul polinomio di grado 9}

\begin{figure}[ht]
\centering
\begin{tikzpicture}
\begin{groupplot}[group style={group size=2 by 1, horizontal sep=0.5cm}, scatola, width=6cm, height=4.2cm,
  xmin=-0.05, xmax=1.05, ymin=-1.6, ymax=1.6, xtick={0,1}, ytick={-1,0,1}, xlabel={$x$}]
\nextgroupplot[title={\small $\ln\lambda = -18$}]
\addplot[thick, verde] table {dati/poly_vera.dat};
\addplot[thick, rosso] table {dati/poly9_reg-18.dat};
\addplot[only marks, mark=o, mark size=1.6pt, blu, thick] table {dati/poly_tr10.dat};
\nextgroupplot[title={\small $\ln\lambda = 0$}, yticklabels={,,}]
\addplot[thick, verde] table {dati/poly_vera.dat};
\addplot[thick, rosso] table {dati/poly9_reg0.dat};
\addplot[only marks, mark=o, mark size=1.6pt, blu, thick] table {dati/poly_tr10.dat};
\end{groupplot}
\end{tikzpicture}
\caption{Polinomio di grado 9 regolarizzato sugli stessi 10 punti della \figrif{fig:polinomi}: con
$\ln\lambda=-18$ la curva è liscia e segue la funzione vera; con $\ln\lambda=0$ è troppo piatta.}
\label{fig:poly-reg}
\end{figure}

Senza regolarizzazione ($\lambda=0$) il polinomio di grado 9 aveva errore di training nullo e overfitting
fortissimo. Con un $\lambda$ adatto ($\ln\lambda = -18$, cioè $\lambda\approx1{,}5\cdot10^{-8}$) il modello è
\textbf{peggiore sul training} ma molto migliore sui dati nuovi. Con $\lambda$ troppo alto ($\ln\lambda=0$,
$\lambda=1$) si finisce in underfitting (\figrif{fig:poly-reg} e \figrif{fig:rms-lambda}). I coefficienti
(Tabella~\ref{tab:coeff-lambda}) mostrano come la penalità ``doma'' i pesi enormi del caso non regolarizzato.

\begin{figure}[ht]
\centering
\begin{tikzpicture}
\begin{axis}[stile, width=8cm, height=4.8cm, xmin=-38, xmax=0.5, ymin=0, ymax=1.05, xlabel={$\ln\lambda$}, ylabel={$E_{RMS}$},
  xtick={-35,-30,-25,-20,-15,-10,-5,0}, ytick={0,0.5,1}, yticklabels={$0$,$0{,}5$,$1$}, legend pos=north west]
\addplot[thick, blu] table[x=lnl, y=train] {dati/poly_rms_lambda.dat}; \addlegendentry{training}
\addplot[thick, rosso, restrict y to domain=0:1.05] table[x=lnl, y=test] {dati/poly_rms_lambda.dat}; \addlegendentry{test}
\end{axis}
\end{tikzpicture}
\caption{$E_{RMS}$ di training e di test del polinomio di grado 9 al variare di $\ln\lambda$: a sinistra
(regolarizzazione debole) overfitting, a destra underfitting, al centro il buon compromesso.}
\label{fig:rms-lambda}
\end{figure}

\begin{table}[ht]
\centering\small
\caption{Coefficienti del polinomio di grado 9 per tre valori di $\lambda$, sugli stessi dati della
\figrif{fig:poly-reg}.}
\label{tab:coeff-lambda}
% generato da script/genera_dati.py: coefficienti del polinomio di grado 9 (lambda = 0, e^-18, 1)
\begin{tabular}{lrrr}
\toprule
 & $\ln\lambda=-\infty$ ($\lambda=0$) & $\ln\lambda=-18$ & $\ln\lambda=0$ \\
\midrule
$w_0^*$ & $0{,}35$ & $0{,}35$ & $0{,}45$ \\
$w_1^*$ & $234{,}52$ & $5{,}44$ & $-0{,}35$ \\
$w_2^*$ & $-5374{,}81$ & $-11{,}96$ & $-0{,}37$ \\
$w_3^*$ & $49046{,}92$ & $19{,}54$ & $-0{,}24$ \\
$w_4^*$ & $-233841{,}49$ & $-82{,}35$ & $-0{,}12$ \\
$w_5^*$ & $645849{,}46$ & $35{,}41$ & $-0{,}02$ \\
$w_6^*$ & $-1070953{,}26$ & $158{,}00$ & $0{,}05$ \\
$w_7^*$ & $1050929{,}32$ & $-19{,}31$ & $0{,}11$ \\
$w_8^*$ & $-562020{,}21$ & $-262{,}48$ & $0{,}16$ \\
$w_9^*$ & $126129{,}45$ & $157{,}62$ & $0{,}20$ \\
\bottomrule
\end{tabular}

\end{table}

\section{Altre forme di regolarizzazione}
\label{sec:l1}

\begin{itemize}
  \item \textbf{Ridge} (norma $\norm{\cdot}_2$): penalizza il quadrato dei pesi e tende a rendere \textbf{tutti} i
        pesi più piccoli.
  \item \termine{Lasso} (norma $\norm{\cdot}_1$, penalità $\lambda\sum_j\abs{w_j}$): tende a portare \textbf{alcuni
        pesi esattamente a zero}, lasciandone altri grandi. È una forma di \textbf{feature selection} e produce
        modelli più semplici. Il valore assoluto rende però la loss non differenziabile in zero. Servono allora
        tecniche di ottimizzazione non differenziabile (sub-gradiente, approssimazioni lisce della norma).
  \item \termine{Elastic net}: usa entrambe le norme.
\end{itemize}

\begin{idea}[Perché $L^1$ azzera i pesi e $L^2$ no]
Il gradiente di $\lambda w^2$ è $2\lambda w$: diventa piccolissimo quando $w$ è vicino a zero. La spinta verso
zero quindi si attenua, e il peso resta piccolo ma non nullo. Il (sub-)gradiente di $\lambda\abs{w}$ è invece
$\lambda\sign(w)$: una spinta \textbf{costante} verso zero, anche per pesi già piccoli. Questa spinta li porta
esattamente a zero se il termine d'errore non li ``difende''. Geometricamente (\figrif{fig:l1-l2}) le curve di
livello della norma $L^1$ hanno \textbf{spigoli sugli assi}, dove una coordinata è nulla. La soluzione del
problema vincolato cade tipicamente su uno spigolo.
\end{idea}

\begin{figure}[ht]
\centering
\begin{tikzpicture}[scale=0.85]
\draw[->, grigio] (-1.6,0) -- (3.6,0) node[right, black, font=\small] {$w_1$};
\draw[->, grigio] (0,-1.4) -- (0,2.9) node[above, black, font=\small] {$w_2$};
\draw[thick, fill=grigiochiaro] (1,0) -- (0,1) -- (-1,0) -- (0,-1) -- cycle;
\draw[rosso, thin] (1.3,1.7) ellipse ({1.044} and {0.4475});
\draw[rosso, thick] (1.3,1.7) ellipse ({2.088} and {0.895});
\fill[rosso] (1.3,1.7) circle (1.4pt) node[right, font=\footnotesize] {$\hat\w$};
\fill[blu] (0,1) circle (2pt) node[left, font=\footnotesize] {$w_1=0$};
\node[font=\small] at (1,-1.85) {$L^1$: $\abs{w_1}+\abs{w_2}\le c$};
\begin{scope}[xshift=5.6cm]
\draw[->, grigio] (-1.6,0) -- (3.6,0) node[right, black, font=\small] {$w_1$};
\draw[->, grigio] (0,-1.4) -- (0,2.9) node[above, black, font=\small] {$w_2$};
\draw[thick, fill=grigiochiaro] (0,0) circle (1);
\draw[rosso, thin] (1.3,1.7) ellipse ({1.002} and {0.4295});
\draw[rosso, thick] (1.3,1.7) ellipse ({2.004} and {0.859});
\fill[rosso] (1.3,1.7) circle (1.4pt) node[right, font=\footnotesize] {$\hat\w$};
\fill[blu] (0.254,0.967) circle (2pt);
\node[font=\small] at (1,-1.85) {$L^2$: $w_1^2+w_2^2\le c$};
\end{scope}
\end{tikzpicture}
\caption{Perché $L^1$ produce soluzioni sparse. $\hat\w$ è il minimo del solo termine d'errore; le ellissi rosse sono
sue curve di livello. La soluzione regolarizzata (blu) è il punto della regione ammessa (grigia) toccato dalla
curva di livello più bassa. Con $L^1$ il contatto avviene su un vertice del rombo, dove $w_1=0$. Con $L^2$ avviene
in un punto generico del cerchio, dove entrambi i pesi sono diversi da zero.}
\label{fig:l1-l2}
\end{figure}

Altri miglioramenti (opzionali):
\begin{itemize}
  \item \termine{data augmentation}: aggiungere rumore agli input insegna al modello a ignorare variazioni
        irrilevanti (robustezza e regolarizzazione). L'\emph{oversampling} aiuta con dataset sbilanciati;
  \item \textbf{input derivati}: si usa un piccolo numero di nuove variabili, combinazioni lineari degli input
        (\ingl{Principal Component Regression}, \ingl{Partial Least Squares}).
\end{itemize}

\begin{puntichiave}
\begin{itemize}
  \item Overfitting: esiste $h'$ con errore di training maggiore ma rischio minore.
  \item Nel fitting polinomiale l'errore di training cala sempre con $M$, quello di test ha forma a U; più dati
        permettono modelli più complessi.
  \item Ridge: $\mathrm{Loss} = \norm{\dd-X\w}^2 + \lambda\norm{\w}^2$; soluzione $(X\T X+\lambda I)^{-1}X\T\dd$;
        weight decay $\w\leftarrow\w+\eta\Delta\w-2\lambda\w$.
  \item $\lambda$ piccolo $\to$ overfitting, $\lambda$ grande $\to$ underfitting; si sceglie con la model
        selection.
  \item $L^1$ azzera alcuni pesi (sparsità), $L^2$ li riduce tutti.
\end{itemize}
\end{puntichiave}

\begin{domande}
\begin{enumerate}
  \item Definire formalmente l'overfitting.
  \item Descrivere l'esempio del fitting polinomiale: cosa succede al variare di $M$ e di $l$?
  \item Perché un errore di training nullo non implica una buona soluzione?
  \item Scrivere la loss di Tikhonov, la soluzione diretta e la regola di weight decay; ruolo di $\lambda$.
  \item Differenza tra regolarizzazione $L^1$ e $L^2$.
\end{enumerate}
\end{domande}
